\subsection{交换环上矩阵的张量积}

设 C 为一个交换环，E、F、U、V 为四个 C-模，且 \(\phi:E\to U\) , \(\psi:F\to V\) 为两个 C-线性映射。假设 E、F、U、V 分别具有有限基 \((e_{\lambda})_{\lambda \in L}, (f_{\mu})_{\mu \in M}, (u_{\rho})_{\rho \in R}, (v_{\sigma})_{\sigma \in S}\) ；令 \(A = (a_{\rho \lambda})\) 为 \(\phi\) 关于 \((e_{\lambda})\) 和 \((u_{\rho})\) 的矩阵， \(B = (b_{\sigma \mu})\) 为 \(\psi\) 关于 \((f_{\mu})\) 和 \((v_{\sigma})\) 的矩阵。对每个有序对 \((\lambda, \mu) \in L \times M = N\) ，令 \(g_{\lambda \mu} = e_{\lambda} \otimes f_{\mu}\) ；对每个有序对 \((\rho, \sigma) \in R \times S = T\) ，令 \(w_{\rho \sigma} = u_{\rho} \otimes v_{\sigma}\) ；于是诸 \(g_{\lambda \mu}\) 构成 \(E \otimes F\) 的一组基，诸 \(w_{\rho \sigma}\) 构成 \(U \otimes V\) 的一组基（ \(\S 3\) ，第 6 号，命题 7 的推论 2）。 \(A\) 与 \(B\) 的张量积，记作 \(A \otimes B\) ，是矩阵 \(X = (x_{\tau v})(\tau, v) \in T \times N\) ，其元素由

\[
x _ {(\rho , \sigma), (\lambda , \mu)} = a _ {\rho \lambda} b _ {\sigma \mu}.\tag{42}
\]

于是 \(A \otimes B\) 是 \(\phi \otimes \psi\) 关于基 \((g_{\lambda \mu})\) 和 \((w_{\rho \sigma})\) 的矩阵。根据定义（§ 3，第 2 号，公式 (3)）

\[
\begin{array}{r l} (\phi \otimes \psi) (g _ {\lambda \mu}) & = (\phi \otimes \psi) (e _ {\lambda} \otimes f _ {\mu}) = \phi (e _ {\lambda}) \otimes \psi (f _ {\mu}) \\ & = \sum_ {\rho , \sigma} a _ {\rho \lambda} b _ {\sigma \mu} (u _ {\rho} \otimes v _ {\sigma}) = \sum_ {\rho , \sigma} a _ {\rho \lambda} b _ {\sigma \mu} w _ {\rho \sigma}. \end{array}
\]

 \(A \otimes B\) 的元素的定义 (42) 表明，这个矩阵与矩阵的矩阵 \((a_{\rho \lambda}B)_{(\rho, \lambda) \in \mathbb{R} \times \mathbb{L}}\) 一一对应，并且也与矩阵 \((Ab_{\sigma \mu})_{(\sigma, \mu) \in \mathbb{S} \times \mathbb{M}}\) 一一对应（第 5 号）。

 \((\phi, \psi) \mapsto \phi \otimes \psi\) 是 C-双线性映射这一事实，连同 §3 第 5 号的公式 (9)，可以用下列恒等式表达：

\[
\left\{ \begin{array}{l} A \otimes (B _ {1} + B _ {2}) = A \otimes B _ {1} + A \otimes B _ {2} \\ (A _ {1} + A _ {2}) \otimes B = A _ {1} \otimes B + A _ {2} \otimes B \end{array} \right.\tag{43}
\]

\[
(c A) \otimes B = A \otimes (c B) = c (A \otimes B) \quad \text { for } c \in \mathbf {C}\tag{44}
\]

\[
(A _ {1} \otimes B _ {1}) (A _ {2} \otimes B _ {2}) = (A _ {1} A _ {2}) \otimes (B _ {1} B _ {2})\tag{45}
\]

当所出现的运算有定义时。矩阵的张量积的转置由下式给出

\[
{ } ^ { t } ( A \otimes B ) = ( { } ^ { t } A ) \otimes ( { } ^ { t } B ) .\tag{46}
\]

如果 A 和 B 是 C 上的可逆方阵， \(A \otimes B\) 是可逆的，并且

\[
(A \otimes B) ^ {- 1} = (A ^ {- 1}) \otimes (B ^ {- 1}).\tag{47}
\]

设 \((e_{\lambda}^{\prime})_{\lambda \in L}\) 是 E 的另一组基， \((f_{\mu}^{\prime})_{\mu \in M}\) 是 F 的另一组基；如果 P 是从基 \((e_{\lambda})\) 到基 \((e_{\lambda}^{\prime})\) 的过渡矩阵，Q 是从基 \((f_{\mu})\) 到基 \((f_{\mu}^{\prime})\) 的过渡矩阵，则从基 \((e_{\lambda} \otimes f_{\mu})\) 到基 \((e_{\lambda}^{\prime} \otimes f_{\mu}^{\prime})\) 的过渡矩阵是 \(P \otimes Q\) 。若 \(A^{\prime}\) 与 A 等价（相应地，相似）且 \(B^{\prime}\) 与 B 等价（相应地，相似），则 \(A^{\prime} \otimes B^{\prime}\) 与 \(A \otimes B\) 等价（相应地，相似）。

矩阵的张量积定义可以以显然的方式推广到C上任意有限多个矩阵；特别地，我们有结合律公式

\[
\left(\bigotimes_ {i \in I _ {1}} X _ {i}\right) \otimes \left(\bigotimes_ {i \in I _ {2}} X _ {i}\right) = \bigotimes_ {i \in I} X _ {i}\tag{48}
\]

对有限指标集 I 的每个划分 \((\mathbf{I}_{1}, \mathbf{I}_{2})\) 成立。

